% Generated draft: requires mathematical and source review.
\documentclass[11pt]{article}
\usepackage[margin=25mm]{geometry}
\usepackage{fontspec}
\setmainfont{Latin Modern Roman}
\usepackage{amsmath,amssymb,mathtools,unicode-math}
\setmathfont{Latin Modern Math}
\usepackage{enumitem,needspace,xurl,hyperref}
\setlength{\parindent}{0pt}
\setlength{\parskip}{5pt}
\setlength{\emergencystretch}{3em}
\newcommand{\N}{\mathbb{N}}
\newcommand{\Z}{\mathbb{Z}}
\newcommand{\Q}{\mathbb{Q}}
\newcommand{\R}{\mathbb{R}}
\newcommand{\C}{\mathbb{C}}
\newcommand{\F}{\mathbb{F}}
\newcommand{\A}{\mathbb{A}}
\newcommand{\PP}{\mathbb{P}}
\newcommand{\E}{\mathbb{E}}
\newcommand{\eps}{\varepsilon}
\newcommand{\dd}{\,\mathrm{d}}
\begin{document}
\section*{UnsolvedMath ID: 30004085}
\textbf{Source problem code:} catalogue entry
\subsection{Definitions and mathematical statement}
For positive integers \(m,k,N\), equip \(\C^m\) and \(\C^N\) with their usual Euclidean topologies. Under the standard complex interpretation, a map \(f:\C^m\to\C^N\) is \emph{\(k\)-regular} if, for every ordered collection of pairwise distinct points \(x_1,\ldots,x_k\in\C^m\), its images are linearly independent over \(\C\); explicitly,
\[
\sum_{i=1}^{k}a_i f(x_i)=0,\qquad a_1,\ldots,a_k\in\C
\quad\Longrightarrow\quad
 a_1=\cdots=a_k=0.
\]
Continuity is required of \(f\); no holomorphic or polynomial condition is imposed in the supplied statement.

The proposed assertion is that, for every admissible triple \((m,k,N)\),
\[
\bigl(\exists\text{ a continuous \(k\)-regular map }f:\C^m\to\C^N\bigr)
\quad\Longleftrightarrow\quad
N\geq m(k-1)+1.
\]
Here the positive-integer parameter range and complex linear independence provide an explicit conventional interpretation: the supplied record does not itself define \(k\)-regularity or specify the parameter range. In particular, whether zero-dimensional spaces or restrictions on \(k\) are intended requires source confirmation; no extension to those unspecified cases is asserted here. According to the supplied verification note, the report presents this assertion as Conjecture 5. It is recorded here as a conjectural equivalence, not as a theorem.
\subsection{Short English statement}
The question asks whether a continuous map from \(\C^m\) to \(\C^N\) can send every \(k\) distinct points to complex-linearly independent vectors exactly when \(N\geq m(k-1)+1\). The definition and parameter conventions require confirmation from the source. The supplied record labels the problem open.
\subsection{Sources}
\begin{itemize}
\item Supplied local UnsolvedMath registry record, permanent ID \(30004085\), title: Optimal Dimensions for Complex \(k\)-Regular Maps; exact source problem code: \texttt{catalogue entry}. The record gives proposed year 2019 and names no individual proposer.
\item Supplied report citation: \emph{Singularities and Homological Aspects of Commutative Algebra} (2020), Oberwolfach Reports, pp.\ 407--464, DOI \texttt{10.4171/owr/2019/6}. \url{https://doi.org/10.4171/owr/2019/6}. The supplied statement-verification note says that the official report labels the assertion Conjecture 5 and that its conversion to a question preserves the threshold. This verification claim has not been independently checked here.
\item Supplied literature source, labeled Buczynski--Januszkiewicz--Jelisiejew--Michalek: \url{https://arxiv.org/abs/1511.05707}.
\item Supplied literature source, labeled Oberwolfach Report 6/2019: \url{https://doi.org/10.4171/owr/2019/6}.
\item Literature assessment dated exactly \texttt{2026-08-22}: the supplied curation labels the problem open and states that the proposed exact threshold is not known in general, with available lower bounds and constructions whose ranges do not universally coincide. This is supplied dated curation derived from the OWR record, not an independently verified literature assessment. The supplied background explicitly calls for independent verification of time-sensitive status and cited claims.
\end{itemize}
\end{document}
